\documentclass{article}
\usepackage{amsmath, amssymb, amsthm}
\usepackage{hyperref}
\usepackage{geometry}
\geometry{margin=1in}

\title{Erd\H{o}s Problem \#508}
\date{Last updated: 2025-08-31}
\author{Source: erdosproblems.com}

\begin{document}
\maketitle

\noindent\textbf{Status:} OPEN \quad \textbf{No prize}

\noindent\textbf{Tags:} geometry, ramsey theory

\medskip
\noindent\textbf{Problem Statement:}

What is the chromatic number of the plane? That is, what is the smallest number of colours required to colour $\mathbb{R}^2$ such that no two points of the same colour are distance $1$ apart?

\bigskip
\noindent\textbf{Remarks:}

The Hadwiger-Nelson problem. Let $\chi$ be the chromatic number of the plane. An equilateral triangle trivially shows that $\chi\geq 3$. There are several small graphs that show $\chi\geq 4$ (in particular the Moser spindle and Golomb graph). The best bounds currently known are\[5 \leq \chi \leq 7.\]The lower bound is due to de Grey \cite{dG18}. The upper bound can be seen by colouring the plane by tesselating by hexagons with diameter slightly less than $1$.

Matolcsi, Ruzsa, Varga, and Zs\'{a}mboki have proved that the fractional chromatic number of the plane is at least $4$. Croft \cite{Cr67} has proved it is at most $4.359\cdots$.

See also [704], [705], and [706]. The independence number of a finite unit distance graph is the topic of [1070].

\bigskip
\noindent\textbf{References:}

[Cr67] H. T. Croft, Incidence incidents. Eureka (1967), 22-26.
 
 [dG18] de Grey, Aubrey D. N. J., The chromatic number of the plane is at least 5. Geombinatorics (2018), 18-31.

\bigskip
\noindent\small{Source: \url{https://www.erdosproblems.com/508}}
\end{document}
